% =============================================================================
%  III. RESEARCH QUESTION AND HYPOTHESES
% =============================================================================
\begin{frame}{III. Research Question and Hypotheses}
\begin{block}{Research Question}
  In real-time sepsis prediction using MIMIC-IV~\cite{johnson2023mimiciv},
  how much of the model's performance difference comes from:
  (a)~the choice of AI model, and (b)~the researcher's decisions, including
  how sepsis is defined, how data is split, when prediction is made, and how
  performance is measured?
\end{block}

\vspace{6pt}
\textbf{Pre-Registered Hypotheses}
\begin{itemize}\itemsep1pt\small
  \item \textbf{H1:} Label choice affects AUROC as much as or more than model
        choice~\cite{cohen2024subtle}.
  \item \textbf{H2:} Feature importance changes across different label
        definitions~\cite{lauritsen2021framing}.
  \item \textbf{H3:} Removing post-treatment information lowers AUROC by
        $\geq$\,0.03~\cite{kamran2024evaluation}.
  \item \textbf{H4:} Temporal testing gives lower AUROC than random testing
        by $\geq$\,0.02~\cite{guo2022temporal}.
  \item \textbf{H5:} Utility Score drops more than AUROC under realistic
        testing~\cite{wang2025srpm}.
  \item \textbf{H6:} Statistical models perform close to gradient boosting
        models (within 0.02 AUROC)~\cite{wang2022equity}.
\end{itemize}
\end{frame}
